\chapter{Hiện thực và kiểm thử}
\label{chap:hienthuc}

% TODO: toàn bộ chương này cần nhóm điền nội dung thực tế từ quá trình phát triển.

\section{Hiện thực}
\label{sec:hienthuc}

% TODO: toàn bộ chương này cần nhóm điền nội dung thực tế từ quá trình phát triển. Khung mục dưới đây bám theo Ch.3 để đảm bảo tính nhất quán thiết kế--hiện thực.

\subsection{Môi trường và công cụ phát triển}
\label{sec:moitruong}
Mã nguồn được quản lý trên GitHub. Mỗi lần đẩy mã, GitHub Actions tự chạy kiểm thử đơn vị kèm đo độ phủ, phân tích tĩnh bằng SonarQube Cloud (bản miễn phí: kho riêng tư tới 50 nghìn dòng mã, kho công khai không giới hạn) và cổng chất lượng; mã mới không đạt ngưỡng bị chặn trước khi hợp nhất. Kiểm thử tải bằng k6 và đo giao diện bằng Lighthouse chạy riêng theo đợt vì tốn thời gian. Nếu quy mô mã nguồn vượt giới hạn của bản miễn phí, chuyển sang SonarQube Community Build tự chạy mà không phải viết lại cấu hình phân tích.

Công nghệ dùng khi hiện thực bám đúng các lựa chọn tại mục~\ref{sec:giaiphapcongnghe}:
\begin{itemize}
    \item \textbf{Kho mã:} một monorepo quản lý bằng pnpm workspaces và Turborepo, gồm ứng dụng mua hàng, ứng dụng nội bộ, ứng dụng lõi và các gói dùng chung (kiểu dữ liệu TypeScript, thư viện thành phần giao diện).
    \item \textbf{Giao diện:} ứng dụng mua hàng bằng Next.js, đóng gói dạng PWA; ứng dụng nội bộ bằng React, Vite và Refine.
    \item \textbf{Ứng dụng lõi:} NestJS trên Node.js và TypeScript; truy cập dữ liệu qua Prisma, riêng khóa dòng tồn viết bằng SQL thô.
    \item \textbf{Dữ liệu:} PostgreSQL kèm các tiện ích pgvector, \texttt{unaccent}, \texttt{pg\_trgm}; Redis cho cache và hàng đợi; MinIO cho tệp.
    \item \textbf{Dịch vụ AI:} Python với FastAPI; mô hình nhúng BGE-M3; thư viện dự báo chuỗi thời gian cho SARIMA, Prophet.
    \item \textbf{Bên thứ ba ở môi trường thử nghiệm:} VNPay, Giao Hàng Nhanh; hóa đơn điện tử giả lập; Web Push (VAPID) và email qua SMTP.
\end{itemize}

% TODO: ghi phiên bản cụ thể của từng thành phần khi dựng xong môi trường; ngưỡng cổng chất lượng đã đặt.

\subsection{Kiến trúc triển khai}
\label{sec:trienkhai}
% TODO: mô tả hạ tầng triển khai thực tế (server/cloud, container hóa nếu có), đối chiếu với kiến trúc thiết kế tại mục~\ref{sec:kientruc}.

\subsection{Hiện thực các module chính}
\label{sec:hienthucmodule}

\subsubsection*{Module xác thực \& quản lý tài khoản}
% TODO

\subsubsection*{Module sản phẩm \& truy xuất nguồn gốc (QR)}
% TODO

\subsubsection*{Module bán lẻ (B2C)}
% TODO

\subsubsection*{Module đánh giá \& uy tín}
% TODO: hiện thực quy trình tại mục~\ref{subsec:danhgianghiepvu} (flowchart Hình~\ref{fig:flow-danh-gia}) -- thu thập đánh giá, kiểm duyệt, tính điểm uy tín tổng hợp.

\subsubsection*{Module bán sỉ (B2B)}
% TODO

\subsubsection*{Module quản lý tồn kho \& vận chuyển/khiếu nại}
% TODO: hiện thực quản lý tồn kho theo lô (FEFO) và tích hợp API với đơn vị vận chuyển 3PL -- mục~\ref{subsec:khobai}.

\subsubsection*{Module quản trị}
% TODO

\subsubsection*{Module AI}
\label{sec:hienthucAI}
% TODO: mô tả hiện thực cụ thể cho 3 mô-đun AI thiết kế tại mục~\ref{sec:thietkeAI} (tìm kiếm theo ngữ nghĩa, gợi ý sản phẩm, dự báo giá) -- mô hình/thư viện sử dụng, dữ liệu huấn luyện, cách tích hợp vào hệ thống chính. Riêng phần giám sát bất thường hiện thực bằng luật ngưỡng, không phải mô hình học máy.

\subsection{Giải thuật và đoạn code quan trọng}
\label{sec:giaithuat}
% TODO: trích các đoạn code/giải thuật tiêu biểu (vd sinh mã QR, tính giá bậc theo MOQ, middleware RBAC, logic escrow) kèm giải thích ngắn gọn.


\section{Kiểm thử}
\label{sec:ketquachucnang}
% TODO: kết quả thực thi các test case thiết kế tại mục~\ref{sec:kichbankiemthu} (đạt/không đạt, lỗi phát hiện và cách khắc phục).

\section{Kết chương}
% TODO
